\section{Physical Model}
\label{sec:physical-model}

\subsection{Scattered-field backends}

Syniscopy computes a complex scattered field by Fourier-propagating a pupil function. In scalar mode this yields a scalar amplitude spread function (ASF); in vectorial Debye mode it yields polarization components $(E_x,E_y,E_z)$. With medium wavenumber $k_0=2\pi n_m/\lambda$, transverse wavevector $\mathbf{k}_\perp=(k_x,k_y)$, pupil amplitude $P(\mathbf{k}_\perp)$, and $k_z(\mathbf{k}_\perp)=\sqrt{k_0^2-|\mathbf{k}_\perp|^2}$ for propagating components,
\begin{equation}
  \mathrm{ASF}(x,y,z)=
  \iint P(\mathbf{k}_\perp)
  \exp\!\left[i(k_xx+k_yy+k_z(\mathbf{k}_\perp)z)\right]
  d^2\mathbf{k}_\perp .
\end{equation}
The normalization is absorbed into $P$ under the chosen Fourier-transform convention. Syniscopy evaluates the field by pupil propagation using the fast Fourier transform \cite{fft-pupil}. For scalar coherent paths, it applies the Lorenz--Mie $S_2(\theta)$ amplitude to the pupil-spreader as a forward-scattering approximation for finite spheres; vectorial Debye applies polarization-resolved pupil fields with $S_1$/$S_2$ terms when material refractive index and diameter are available. Syniscopy caches the result as a stack of two-dimensional complex fields over axial position for each unique sub-particle type. The default optical backend is vectorial Debye (polarization-aware pupil and field-component handling); scalar paraxial remains available through the \texttt{scalar\_paraxial} selector.

\subsection{Imaging models}

Let $E_{\mathrm{sca}}(\mathbf{r},t)$ denote the complex scattered field at the image plane. The observed detector-domain quantity depends on how that scattered field combines with a reference or incident field:
\begin{align}
I_{\mathrm{DF}} &= |E_{\mathrm{sca}}|^2, &
I_{\mathrm{COBRI}} &= |E_{\mathrm{inc}}+E_{\mathrm{sca}}|^2, \\
\phi_{\mathrm{QPI}} &= \arg\!\left(1+\frac{E_{\mathrm{sca}}}{E_{\mathrm{ref}}}\right), &
I_{\mathrm{RICM}} &= |\rho_sE_{\mathrm{ref}}+\rho_p e^{i\phi_{\mathrm{int}}}E_{\mathrm{sca}}|^2, \\
I_{\mathrm{DHM}} &= |E_{\mathrm{ref}}e^{i\mathbf{K}\cdot\mathbf{r}}+E_{\mathrm{sca}}|^2, \\
I_{\mathrm{iSCAT}} &= |E_{\mathrm{ref}}+E_{\mathrm{sca}}|^2 .
\end{align}
Here dark-field microscopy, coherent bright-field imaging, quantitative phase imaging, reflection interference contrast microscopy, digital holographic microscopy, and interferometric scattering microscopy correspond to the displayed detector-domain observables. $E_{\mathrm{inc}}$ denotes the incident bright-field component, $E_{\mathrm{ref}}$ the explicit reference field, $\rho_s$ and $\rho_p$ substrate and particle-interface reflection coefficients, $\phi_{\mathrm{int}}$ the interface phase accumulated by the reflected/scattered components, and $\mathbf{K}$ the off-axis holography carrier wavevector.

The coherent-optical label is used here in the scattered-field-composition sense, with backend selection controlling the assumptions. Syniscopy exposes coherent bright-field imaging and a partially coherent K\"{o}hler-style bright-field path. It likewise exposes annular K\"{o}hler dark-field as a primary dark-field path and coherent dark-field as a reference-free coherent limit. Cross-modality coherent diagnostics use either scalar or vectorial coherent paths according to \texttt{optical\_field\_backend} configuration, while keeping consistent detector-domain observables.

Zernike phase contrast and differential phase contrast are represented as phase-ring and asymmetric-illumination families, respectively. Zernike phase contrast defaults to an objective-pupil phase-ring transfer normalized by numerical aperture and wavelength. Differential-phase contrast defaults to a vectorial Debye half-pupil asymmetric-illumination intensity transfer and supports \texttt{analyzer\_x}, \texttt{analyzer\_y}, \texttt{incoherent\_sum}, \texttt{unpolarized}, and \texttt{full\_vector} detection modes, with \texttt{full\_vector} as the default vectorial observable. The phase-gradient and scalar two-axis options remain available only through explicit lower-fidelity backend settings. Widefield fluorescence and total internal reflection fluorescence select between a parametric emission point-spread-function profile and the vectorial photophysics backend.

Transmission electron microscopy now uses \texttt{multislice\_physical} as the comparison-profile backend. The backend constructs a material-derived projected-potential source, applies split-step Fresnel propagation through configured slices, applies objective-aperture and contrast-transfer terms in Fourier space, and records whether an external reference-backend validation hash is present. The lightweight \texttt{ctf\_proxy}, \texttt{multislice\_lite}, and \texttt{syniscopy\_multislice} routes remain explicit diagnostic backends rather than the paper comparison default.

Scanning electron microscopy now uses \texttt{physical\_electron\_transport} with \texttt{monte\_carlo\_physical} as the comparison-profile backend. The renderer preserves material-resolved SEM source channels, so the transport backend convolves each material channel with a material-specific kernel rather than collapsing gold, carbon, water, and support material into one numeric source map. The current production comparison uses screened-Rutherford elastic scattering with Joy--Luo continuous stopping and energy-cutoff transport; a Browning-style empirical Mott surrogate is exposed through \texttt{sem\_physical\_elastic\_model=\allowbreak mott\_browning} for explicit sensitivity studies \cite{joy-luo-stopping,browning-mott-cross-sections}. The screened-Rutherford branch is retained as the baseline because the Browning surrogate, inside this simplified transport geometry, over-predicts the empirical 20~keV backscatter-coefficient anchor used by the current validation script \cite{tabata-backscatter}. The SEM validation script also compares material-dependent ranges against the Kanaya--Okayama electron range law \cite{kanaya-okayama-range}. Legacy raster transport, Gaussian-probe proxy, interaction-volume proxy, and reference-kernel table modes remain selectable only as named lower-fidelity or externally calibrated paths. Each implemented modality profile exposes the same renderer hooks---illumination field, response function, sample-environment interaction, and noise model---and records the selected backend fidelity in profile-card metadata.

\subsection{Composite rigid particles}

Rigid composite particles are assemblies of $N$ sub-particles with body-fixed offsets $\mathbf{o}_i$ from the composite center of mass. For center-of-mass position $\mathbf{r}_{\mathrm{cm}}(t)$ and orientation quaternion $q_t$,
\begin{equation}
  \mathbf{r}_i(t)=\mathbf{r}_{\mathrm{cm}}(t)+R(q_t)\mathbf{o}_i .
\end{equation}
The composite scattered field is the coherent sum of cached per-type axial field slices evaluated at each lab-frame sub-particle position. Composite assemblies are therefore linear superpositions of independently propagated sub-particle fields; the current implementation does not model near-field electromagnetic coupling, plasmon hybridization, or multiple scattering between closely spaced components unless those effects are absorbed into a user-supplied effective scatterer model. Because scalar pupil propagation is cached per sub-particle type, each frame places $N$ transformed component fields rather than rerunning the pupil propagation for every component. The generated component-work record is therefore an operation-count diagnostic for cached field placement, not a wall-clock benchmark.

Translational trajectories are Brownian. Rotational trajectories are Brownian on the three-dimensional rotation group $SO(3)$ with either a user-specified angular step or an opt-in Stokes--Einstein/Perrin-style estimate from hydrodynamic diameter, viscosity, and temperature \cite{stokes-einstein,rotational-diffusion}. These dynamics are configurable synthetic trajectory priors; surface pinning, hindered diffusion, flow, drift, hydrodynamic wall corrections, and particle-substrate adhesion are not implied unless supplied by the user's trajectory model.

\subsection{Sample environment}
\label{sec:substrate}

Syniscopy splits the rendered scene into particle and sample environment. The sample environment contains a surrounding medium, an optional mounting interface, and an optional pattern overlay. The medium supplies refractive index and viscosity for optical contrast and Brownian motion. When present, the mounting interface supplies height and material maps; suspension-only configurations omit that interface. The current implementation includes configurable patterned environments such as gold-hole arrays and nanopillar-style patterns. Additional patterned-support categories (for example, fiducial arrays, grid bars, holey-carbon supports, or microfluidic walls) require explicit model extension or user-supplied pattern maps before they should be treated as implemented presets.

Each modality maps the shared environment into its own contrast channel. This particle/environment separation is what makes substrate-versus-particle ablation studies possible without re-architecting a modality implementation: the same particle can be swept across substrates, and the same substrate can be reused across particle families. Bright-field-style paths read transmission modulation. Fluorescence reads substrate autofluorescence and excitation modulation. Dark-field reads pattern-edge scattering through a high-pass response.

Phase and interference profiles consume the environment through their own reference or phase channels. Quantitative phase imaging, reflection interference contrast microscopy, digital holographic microscopy, and interferometric scattering microscopy read thin-film reflection/transmission into the relevant reference field, while coherent bright-field imaging remains available as a coherent bright-field limit. Transmission electron microscopy adds the mounting-interface projected potential before the contrast-transfer-function response. Scanning electron microscopy adds material-dependent secondary-electron yield and topographic edge terms. Real interferometric scattering recordings can additionally exhibit coherent speckle from substrate inhomogeneity; Syniscopy now includes an optional configurable roughness/speckle perturbation model (\texttt{sample\_environment\_pattern\_roughness\_*}) with static, flicker, or source-matched modes, tunable correlation scale, and phase jitter. Source terms can be coupled through \texttt{independent}, \texttt{coherent\_amplitude}, \texttt{field\_weighted}, \texttt{scene\_weighted}, or \texttt{channel\_weighted} roughness modes (default: \texttt{channel\_weighted}).

\subsection{Noise and supervision model}
\label{sec:supervision}

Intensity and fringe modalities use a Poisson--Gaussian count model after the imaging model maps intensity to detected counts: per-pixel mean $\mu(\mathbf{r})$ has variance $\mu(\mathbf{r})+\sigma_r^2$. The same interface accepts per-modality counts-domain defaults, including per-pixel gain mapping, read-noise counts, dark offset, fixed-pattern gain and offset terms, hot pixels, and scanning electron microscopy scan-line noise. Effective electron gain and excess-noise behavior are represented by public configuration controls such as \texttt{detector\_qe}, \texttt{emccd\_excess\_noise\_factor}, and read-noise maps. A real flat-field recording can also be used to estimate an effective camera gain in electrons per count by inverting the shot-noise relation between temporal variance and mean counts. Drift, vibration, empirical shading, and detector-defect terms are exposed as configurable correction channels (gain/offset maps, scan-line, and field-dependent defect terms) that can be composed through the detector-domain pipeline; higher-order couplings can be approximated by additional source/environment channels when configured. Bench-specific absolute calibration requires calibrated instrument parameters, and Appendix~\ref{sec:bench-mapping} maps common bench measurements to the corresponding configuration parameters. Quantitative phase imaging reports phase in radians and uses a phase-noise model.
\begin{equation}
  \sigma_\phi^2(\mathbf{r}) = \frac{1}{V^2 n_Q} + \sigma_{\phi,r}^2 ,
\end{equation}
where $V$ is visibility, $n_Q$ is per-pixel detected-quanta budget, and $\sigma_{\phi,r}^2$ is optional additive phase-readout variance.

For each simulated particle frame, the model defines a geometry/support mask, a supported mask, an ignore mask, an optional loss-weight map, and an optional rejection label. The geometry/support mask contains the projected object footprint together with the configured contrast-lobe support before supervision gating. Syniscopy produces the supported mask from four bounded support factors. Signal support comes from frame-local contrast-to-noise; information support comes from the Cram\'{e}r--Rao lower bound; temporal support comes from the Brownian step plausibility score; and assignment-ambiguity support comes from overlap or proximity to competing particles. Signal, information, and temporal support are particle-frame factors applied to known-object pixels; ambiguity can vary spatially where competitor masks overlap. The ambiguity factor defaults to one when no competing-particle state is supplied, which is the uniform-assignment special case. If competitor positions or masks are supplied, overlap and proximity reduce the factor and can route known pixels to the ignore mask.

The schema is intended for cases where per-frame contrast is too low or too entangled with structured background for classical thresholding methods such as Otsu, blob detection, or band-pass filters to produce reliable particle masks. The schema is segmenter-agnostic: ignore regions, supported masks, and per-pixel loss weights are standard tensors that Segment Anything Model 2, Cellpose, custom U-Net models, or point trackers can use. Pixels where the simulator knows an object exists but the selected support policy fails are routed to the ignore mask, preventing unsupported known-object pixels from becoming background negatives during training.

\paragraph{Supported supervision as reject labeling.}
Operationally, the rule prevents object pixels known to the simulator from becoming background negatives when the rendered measurement does not support a reliable positive label. The support policy is motivated by Bayes' rule with a reject option \cite{chow-reject}, but the implemented decision is a bounded support-score rule. For known-object pixels in a particle frame, the positive-supervision decision compares an implemented support score against a threshold; pixels below threshold are rejected into the ignore mask rather than relabeled as background.

\paragraph{Support-score decomposition.}
Under the simulator's generative assumptions (Brownian latent dynamics, known particle geometry, detector-domain Gaussian or locally Gaussian noise, and a particle-presence prior), an idealized posterior-odds argument suggests image-likelihood, state-uncertainty, temporal-prior, and assignment terms. Syniscopy evaluates bounded proxies for those terms. For particle $i$ in frame $t$, let $C_i(\mathbf{r},t)$ be its per-particle contrast image, $\sigma_n$ the estimated contrast-noise standard deviation, $\sigma_{xy}$ the lateral Cram\'{e}r--Rao bound, $\sigma_{\max}$ the configured acceptable bound (defaulting to one pixel), $\Delta\mathbf{x}_{i,t}$ the lateral Brownian step, $D_i$ the Stokes--Einstein diffusion coefficient used by the synthetic trajectory prior, $\Delta t$ the frame interval, and $G_i$ the geometry/support mask. The particle-frame support factors are
\begin{equation}
  s_{\mathrm{signal}}
  =
  1-\exp\!\left[-\frac{1}{2}
  \left(\frac{\max_{\mathbf{r}} |C_i(\mathbf{r},t)|}{\sigma_n}\right)^2\right],
  \qquad
  s_{\mathrm{information}}
  =
  \frac{1}{1+(\sigma_{xy}/\sigma_{\max})^2},
\end{equation}
with invalid, non-finite, or singular information mapped to zero support. The first frame of a particle track has $s_{\mathrm{temporal}}=1$; subsequent frames use
\begin{equation}
  s_{\mathrm{temporal}}
  =
  \exp\!\left[
  -\frac{1}{2}
  \left(
  \frac{\lVert\Delta\mathbf{x}_{i,t}\rVert}
  {3\sqrt{4D_i\Delta t}}
  \right)^2
  \right],
\end{equation}
clipped to $[0,1]$. The factor of three is the implemented permissive radial envelope: one two-dimensional Brownian step has radial scale $\sqrt{4D_i\Delta t}$, and the support score decays after roughly three such radial scales rather than at the one-scale displacement. Let $a$ be the configured ambiguity-distance scale and let $\mathcal{C}$ be the set of competing particles. The frame-level ambiguity competition score is
\begin{equation}
  q_i =
  \sum_{j\in\mathcal{C}}
  \exp\!\left[-\frac{1}{2}
  \left(\frac{\lVert\mathbf{x}_{j,t}-\mathbf{x}_{i,t}\rVert}{a}\right)^2\right]
  +
  \sum_{j\in\mathcal{C}}
  \frac{|G_i\cap G_j|}{|G_i|},
  \qquad
  s_{\mathrm{ambiguity}}^{0}=\frac{1}{1+q_i}.
\end{equation}
When no competitor positions or masks are supplied, $q_i=0$ and $s_{\mathrm{ambiguity}}^{0}=1$. If overlap masks are supplied, the pixelwise ambiguity factor is additionally bounded by local overlap count $n_{\mathrm{overlap}}(\mathbf{r})$:
\begin{equation}
  s_{\mathrm{ambiguity}}(\mathbf{r})
  =
  \min\!\left(
  s_{\mathrm{ambiguity}}^{0},
  \frac{1}{1+n_{\mathrm{overlap}}(\mathbf{r})}
  \right).
\end{equation}
Thus signal, information, and temporal terms are particle-frame scores, while overlap ambiguity can vary within the object mask.

The default decision rule converts each included bounded factor to clipped log odds and adds a prior log-odds offset:
\begin{equation}
  L(\mathbf{r})
  =
  L_0
  +
  \sum_{k\in\mathcal{J}}
  \log
  \frac{\operatorname{clip}_{\epsilon}(s_k(\mathbf{r}))}
       {1-\operatorname{clip}_{\epsilon}(s_k(\mathbf{r}))},
  \qquad
  \mathcal{J}\subseteq
  \{\mathrm{signal},\mathrm{information},\mathrm{temporal},\mathrm{ambiguity}\}.
\end{equation}
Known-object pixels are emitted in the supported mask when $L(\mathbf{r})$ exceeds the configured log-odds threshold. An explicit product rule is also available, in which the emitted supported mask contains known-object pixels whose product of included support factors exceeds the configured product threshold.

The default public configuration enables all four factors, uses $\epsilon=10^{-12}$ for clipped log odds, sets $L_0=0$ and the log-odds threshold to zero, and keeps the product-rule threshold at 0.2 for runs that opt into product scoring. The default acceptable information scale $\sigma_{\max}$ is one detector pixel unless a run supplies a tighter localization threshold.

\emph{Motivation.} Bayes' rule writes ideal correct-label posterior odds as a likelihood ratio times prior odds. The detector likelihood motivates the signal term. A local Gaussian/Laplace view of the latent particle-state posterior motivates a Fisher-information uncertainty term \cite{bishop-prml}. Brownian motion motivates the transition-support term between adjacent frames. Marginalizing the discrete assignment of a pixel among candidate particles motivates the ambiguity term. The particle-presence prior is absorbed into the decision threshold; the four support factors above are the quantities evaluated by the implementation.

This is a support-score construction rather than a calibrated posterior identity. The implementation exposes the corresponding bounded factors and clipped log-odds components so the ignore-mask decision can be audited. With no competitor state, the ambiguity factor equals its no-competitor default, so the policy specializes to signal, information, temporal, and default ambiguity terms.
